\documentclass[varwidth=16cm, border=2mm]{standalone}
\usepackage{amsmath}\usepackage{amssymb}\usepackage{ctex}
\usepackage{tasks}\usepackage{xcolor}\usepackage{graphicx}
\settasks{label=\Alph*., label-width=1.5em, item-indent=2em}
\begin{document}
\textbf{[0.6]} 如图所示，由半椭圆 $C_1:\frac{x^2}{4}+\frac{y^2}{b^2}=1(y\leq0)$ 和两个半圆 $C_2:(x+1)^2+y^2=1(y\geq0)$、$C_3:(x-1)^2+y^2=1(y\geq0)$ 组成曲线 $C:F(x,y)=0$，其中点 $A_1,A_2$ 依次为 $C_1$ 的左、右顶点，点 $B$ 为 $C_1$ 的下顶点，点 $F_1,F_2$ 依次为 $C_1$ 的左、右焦点。若点 $F_1,F_2$ 分别为曲线 $C_2,C_3$ 的圆心。(1) 求 $C_1$ 的方程;

\begin{center}\includegraphics[width=0.45\textwidth]{E:/001_Math/figures/q_images/fig_D067556E.png}\end{center}

\vspace{0.2cm}\par\noindent{\color{red}\textbf{【答案】} $C_1:\frac{x^2}{4}+\frac{y^2}{3}=1(y\leq0)$}
\par\noindent{\color{blue}\textbf{【解析】} (1) 由两圆的方程知：圆心分别为 $C_1(-1,0), C_2(1,0)$，即 $F_1(-1,0), F_2(1,0)$，$\therefore b^2+1=4$，解得：$b^2=3$。$\therefore C_1:\frac{x^2}{4}+\frac{y^2}{3}=1(y\leq0)$。}


\end{document}
